\section{Hecke matrix coefficients and support detection}
\label{meas:section}

Fix an abstract coefficient isomorphism $\E\simeq\C$ and a closed point
$v\in X$ of degree $d$.  Put $r=q^d$.  We use the compact sets $B_o$ and
the bilinear pairing $[-,-]$ of \eqref{intro:compactsets} and
\eqref{intro:pairing}.  Write
\[
 B=\bigcup_{o\subset\Nhat}B_o,
 \qquad B_Y=\bigcup_{o\subset Y}B_o.
\]
There are finitely many ambient nilpotent orbits.  Lemma
\ref{desc:compactsets}, whose proof is independent of this section, says
that the $B_o$ are disjoint compact sets and that representation characters
are uniformly dense in $C(B)$.  We will use this elementary fact for
uniqueness and for separating different orbit contributions.

\begin{theorem}[Bilinear support detector]\label{meas:detector}
For every $f,g\in\cA_{\C}$ there is a unique finite complex Borel measure
$\nu_{f,g}$ on $B$ such that, for every finite-dimensional representation
$D$ of $\Ghat$,
\begin{equation}\label{meas:moments}
 [T_Df,g]=\int_B\chi_D\,d\nu_{f,g}.
\end{equation}
For every closed invariant subset $Y\subset\Nhat$,
\begin{equation}\label{meas:criterion}
 f\in\cF_Y\cA_{\C}
 \quad\Longleftrightarrow\quad
 \operatorname{supp}(\nu_{f,g})\subset B_Y
 \quad\text{for every }g\in\cA_{\C}.
\end{equation}
\end{theorem}

The theorem concerns ordinary finitely supported functions.  In its proof,
infinite sums occur only in auxiliary character calculations on compact
groups.  Every automorphic input and test remains a finite linear
combination of compact Weil classes.

\subsection{The bilinear Hom formula}

We first identify exactly which scalar trace is being calculated.  Use the
nonstandard self-duality and compactness of Theorem~\ref{inp:duality}
\cite[Theorem~3.2.2 and Appendix~A.1]{AGKRRV2}.
On the automorphic side its evaluation is
\[
 (\mathcal F,\mathcal G)\longmapsto
 \Gamma_c\bigl(\Bun_G,\mathcal F\mathbin{\overset{*}{\otimes}}
                         \mathcal G\bigr).
\]
On compact objects the categorical dual is the opposite category.
Simultaneous invariance under $\Phi$ identifies its action there with
$\Phi^{\mathrm{op}}$.  Thus, on the spectral side, take a covariant class
$b:N\longrightarrow\Phi N$ and an opposite-input class
$a:\Phi M\longrightarrow M$.  Evaluation acts on an arrow $h:M\to N$ by
\begin{equation}\label{meas:hom-action}
 S_{M,N}(h)=\Phi^{-1}(b\circ h\circ a).
\end{equation}
The action with a Hecke insertion includes its specified cyclic Weil
transport.  Formula~\eqref{meas:hom-action} is the inverse transported action
$S$ of Section~\ref{mod:section}; it uses the original $\Phi$-trace.

\begin{lemma}\label{meas:hom-functions}
Let $N$ and $M$ be compact spectral objects with the structures just
specified, and let $f_N$ and $g_M$ be their automorphic trace functions,
where $M$ is transported through the categorical duality.  Then
\begin{equation}\label{meas:hom-function-formula}
 [T_Df_N,g_M]
 =\operatorname{str}\!\left(
 S_{M,N,D};\operatorname{RHom}(M,\mathcal E_{D,v}\otimes N)\right).
\end{equation}
Here $\mathcal E_{D,v}$ denotes the evaluation insertion at the cyclic tuple
above $v$.  The scalar evaluation complex is perfect.  Finite linear
combinations of these classes give all inputs and all tests in
$\cA_{\C}$.  Inputs in $\cF_Y\cA_{\C}$ may be represented by such $N$ with
the prescribed $Y$-support.
\end{lemma}

\begin{proof}
The compact embedding into ambient sheaves, preservation of compactness
by the star tensor product in two compact variables, and preservation by
the finite Hecke insertion imply that the displayed compact-cohomology
evaluation is perfect.  Compact pushforward to a point preserves
compactness; its target is the category of complexes of vector spaces,
whose compact objects are precisely the perfect complexes.  These are the
compactness assertions recorded with the duality input in
Section~\ref{inp:section}.  They are needed in addition to abstract
self-duality.

The ordinary local-term comparison of Theorem~\ref{inp:trace} and the
sheaf--function dictionary
identify the Frobenius trace of this evaluation with
\[
 \sum_{b\in\Bun_G(\mathbf F_q)/\cong}
       \frac{(T_Df_N)(b)g_M(b)}{|\Aut(b)|}.
\]
For each fixed Hecke label the correspondences are bounded, and the compact
inputs have quasi-compact $!$-support.  Thus the compact local-term theorem
applies to precisely this evaluation.  Transport through the duality gives
\eqref{meas:hom-action}, proving
\eqref{meas:hom-function-formula}.  Finally, Theorem~\ref{tr:spanning}
gives actual compact Weil-class spanning on both the covariant and opposite
sides, and on every closed-support trace image.  Combined with the
trace--functions isomorphism, it proves the last assertions.
\end{proof}

The duality need not preserve individual spectral components or their
support labels.  We choose $M$ on the Hom side and then obtain its function
test by duality.  Distinct components are orthogonal because their spectral
Hom complexes vanish.

\subsection{Supported traces on an orbit}

Fix a retained Frobenius-stable component and the model of
Theorem~\ref{mod:model}.  In the curved description, write
$W(x,z)=\langle z,\mu(x)\rangle$.  For two compact lifts let
$\mathcal H=\mathcal H\!om(M,N)$ be their differential Hom, computed with
locally free representatives.  Its differential squares to zero.
Its derived $J$-invariant sections, with the evaluation insertion, compute
the right side of \eqref{meas:hom-function-formula}.  The support of
$\mathcal H$ lies in the intersection of the two factorization supports:
where either object vanishes in the local factorization category, the
differential Hom is acyclic.

We filter these sections by \emph{local cohomology with supports} in the
nilpotent $z$-orbits.  Ordinary restriction to an orbit would give the wrong
answer on boundary strata.  The following calculation fixes both the
normal-direction sign and the convergence issue.

For the motivating use of local cohomology, normal symmetric powers and
distributional convergence, see Kazhdan--Okounkov
\cite[Sections~2.6.1--2.6.4, equations~(107)--(111)]{KazhdanOkounkov}.
The following calculation supplies the supported convergence needed for
the present bilinear detector.

\begin{lemma}[Supported convergence]\label{meas:supported-convergence}
Let $O=J.e$ be a stable nilpotent orbit, closed in an invariant open subset
of $\mathfrak j$.  The contribution with support on $O$ to the Hom
calculation has an absolutely convergent alternating character on the
compact transporter coset of $L_e$.  Convergence is uniform on that coset,
also after any fixed finite-dimensional evaluation insertion.  These
characters are additive through the finite orbit filtration and compute
the ordinary supertrace of the final perfect evaluation complex.
\end{lemma}

\begin{proof}
Let $i:O\hookrightarrow U\subset\mathfrak j$ be the closed immersion and
$c=\operatorname{codim}_{U}(O)$.  It is a regular immersion, with normal
bundle whose fiber at $e$ is $N_e=\mathfrak j/[\mathfrak j,e]$.  If $I$ is
its ideal, local cohomology is the filtered colimit
\[
 R\Gamma_O(\mathcal H)
 =\underset{n}{\operatorname{colim}}\,
 R\mathcal H\!om(\mathcal O_U/I^{n+1},\mathcal H).
\]
The successive derived cofibers of this filtration are
\begin{equation}\label{meas:normal-filtration}
 i_*\left(Li^*\mathcal H\otimes\Sym^n N
                         \otimes\det N[-c]\right),\qquad n\geq0.
\end{equation}
Indeed, $I^n/I^{n+1}=\Sym^n N^*$, while regular-immersion duality gives
$i^!\mathcal H=Li^*\mathcal H\otimes\det N[-c]$.  One may verify the same
formula directly by the regular Koszul resolution.  Each such resolution
is bounded before folding the grading, so it also applies to the locally
free differential Hom used here.  In particular the determinant in
\eqref{meas:normal-filtration} is the \emph{normal} determinant.

For example, a degree-two coordinate $z$ with $S(z)=qz$ gives
\begin{equation}\label{meas:one-normal-coordinate}
 H^1_{(z)}\C[z]=\bigoplus_{m\geq1}\C z^{-m},
 \qquad S(z^{-m})=q^{-m}z^{-m}.
\end{equation}
The total degree of $z^{-m}$ after including local-cohomology degree is
$1-2m$.  This checks the normal sign and the shift simultaneously.

Take the fiber at $e$ and then derived $C_e=L_e\ltimes U_e$ invariants.
The cohomology of the restricted differential Hom is a coherent graded
module over $\C[V]$.  Its support is contained in $\ker e$.  To see the
last assertion directly, differentiate the factorization differential:
if $d_N^2=W$, then
\[
 d_N(\partial d_N)+(\partial d_N)d_N=\partial W.
\]
Left composition with $\partial d_N$ is a homotopy for multiplication by
$\partial W$ on differential Hom.  At $z=e$ the derivatives in the
$V$-directions give the linear equations $ex=0$.  This also proves the
assertion using any finite locally free replacement.

Coherence implies that some power of the ideal of $\ker e$ annihilates
this cohomology.  Consequently only \emph{finitely many} polynomial orders
transverse to $\ker e$ in the $x$-space occur.  The orders $n$ in
\eqref{meas:normal-filtration}, in contrast, remain unbounded.
After the triple translation, the weights from
Proposition~\ref{orb:hom-calculation} give
\[
 \begin{array}{c|c|c}
 \text{space}&\text{absolute }S\text{-weight}&\text{bound}\\ \hline
 (\ker e\cap V_i)^*&-1-i\quad(i\geq0)&\leq-1\\
 (N_e)_i&-2+i\quad(i\leq0)&\leq-2\\
 \Lie(U_e)_i^*&-i\quad(i>0)&<0.
 \end{array}
\]
Absolute weight $w$ means eigenvalue modulus $q^{w/2}$.  Thus the
unrestricted $x$-variables and the normal symmetric powers both contract.
The finitely many transverse $x$-orders, the normal determinant, and the
finite complex of $U_e$-cochains affect only a finite coefficient factor.

Here is a precise summability bound.  Choose a finite-dimensional invariant
generating space for the restricted cohomology, including its finite
transverse thickening.  The locally finite models constructed in the proof of
Proposition~\ref{tr:orbit-trace} permit this choice for
the algebraic group generated by the normalized transporter and
equivariance.  There are constants $C,a,b$ and $0<\rho_x,\rho_z<1$ such
that the sum of eigenvalue moduli, counted with multiplicities, in the
terms of $x$-degree $m$ and normal order $n$ is at most
\begin{equation}\label{meas:geometric-bound}
 C(m+1)^a(n+1)^b\rho_x^m\rho_z^n.
\end{equation}
For example one can take $\rho_x=q^{-1/2}$ and $\rho_z=q^{-1}$ after
absorbing the finite coefficient spaces into $C$.  Polynomial dimension
growth follows from the finite number of polynomial variables.  Subquotients
can only decrease the multiplicity bound on each generalized weight.

By Lemma~\ref{mod:compact-normalization} these estimates are uniform when a
compact element of $L_e$ is included in the transporter.  The semisimple
normalized parts preserve a compact form; the residual unipotent commutes
with it.  Hence compact twists do not change the contracting moduli of the
linear data.  Finite coefficient representations supply only a uniform
constant.  This reasoning concerns traces and eigenvalue multiplicities;
it does not require a unipotent operator to be unitary.

For completeness, the bound also justifies passage from the graded pieces
to their colimit.  For every $\epsilon>0$, only finitely many pairs $(m,n)$
can have eigenvalues of modulus at least $\epsilon$, since the coefficient
weights lie in a fixed finite set.  In the category of locally finite
representations, taking a sum of generalized eigenspaces in such a range
is exact.  On that range the filtered calculation therefore stabilizes
after finitely many normal orders, and the ordinary long exact sequences
give trace additivity.  Filtered colimits of vector spaces are exact, so
this identifies the same part of the cohomology of the colimit.  Finally
\eqref{meas:geometric-bound} is summable in $(m,n)$, allowing
$\epsilon$ to decrease to zero.  The argument applies also to the finite
orbit filtration.  It uses direct-sum folding of the graded complexes;
bounded Koszul and group-cochain calculations commute with this folding.

If Frobenius permutes a finite set of strata without fixing any, use a
power $S^t$ to make the preceding estimates on each summand.  The same
geometric bound is summable after raising each eigenvalue modulus to
any positive power, in particular $1/t$.  It therefore gives summability
also for $S$, whose one-step action has zero trace on the resulting
permuted sum.  On fixed strata the
preceding calculation applies directly.  The full scalar evaluation is
perfect by Lemma~\ref{meas:hom-functions}, so its convergent alternating
character is its ordinary finite supertrace.
\end{proof}

\subsection{Constructing the measures}

Let $L_{e,c}$ be the maximal compact subgroup of $L_e$ chosen in
Lemma~\ref{mod:compact-normalization}, including all components.  Denote by
$\tau_e$ the triple-adjusted $S$-action, using that normalized transporter
and retaining the absolute-weight scalars in the displayed table above.
On coefficient modules in lifted degree zero it is the normalized
transporter itself.  Expressions on the corresponding compact coset are
written as functions of $k\in L_{e,c}$, with operator $k\tau_e$ on the
space in question.  Haar measure $dk$ is probability measure on the entire
compact group.  This notation allows $\tau_e$ to act nontrivially on $L_e$.

Projection to reductive invariants is integration over $L_{e,c}$:
\begin{equation}\label{meas:compact-projection}
 \operatorname{str}(\tau_e;H^{L_e})
 =\int_{L_{e,c}}\operatorname{str}(k\tau_e;H)\,dk.
\end{equation}
For a finite-dimensional representation this is the Reynolds projection,
which commutes with $\tau_e$ because it normalizes the compact group.
Lemma~\ref{meas:supported-convergence} proves the formula for all the
convergent characters here by uniform absolute convergence.  It also
applies when $L_e$ is disconnected: averaging over its finitely many
components is part of the same Reynolds projection.

The Hecke insertion contributes a finite-dimensional cyclic trace.  For
linear maps around $d$ copies of $D$ one has
\begin{equation}\label{meas:cyclic-trace}
 \operatorname{tr}\bigl(\operatorname{cyc}\circ
                      (A_1\otimes\cdots\otimes A_d);D^{\otimes d}\bigr)
 =\operatorname{tr}(A_d\cdots A_1;D).
\end{equation}
Expanding in a basis proves the identity: the cyclic identification of
indices leaves precisely the matrix entries of the indicated product.
The framed evaluation description in Theorem~\ref{mod:model} applies
this identity to the geometric points above $v$.  The triple translations
contribute $\lambda_e(q^{d/2})$, and the residual semisimple product is
compact and centralizes the triple.  Unipotent factors may be discarded
in this trace.  We obtain a continuous map
\begin{equation}\label{meas:parameter-map}
 \beta_e:L_{e,c}\longrightarrow B_{\Ghat.e}.
\end{equation}
Any lifted-degree sign on the insertion contributes
$\lambda_e(-1)$ to the compact part.  This is already included in
$B_{\Ghat.e}$.  Contravariant character conventions give the inverse
parameter, which lies in the same compact set by
Lemma~\ref{desc:compactsets}.

Remove the insertion and write $a_{M,N,e}(k)$ for the alternating character
density of the supported contribution.  The lifted-degree sign may be
used inside \eqref{meas:compact-projection}: after taking invariants it is
the true cohomological sign, as established in
Proposition~\ref{orb:hom-calculation}.  The density is continuous and
absolutely integrable by Lemma~\ref{meas:supported-convergence}.
The contribution of $O$ is therefore the finite complex measure
\begin{equation}\label{meas:orbit-measure}
 (\beta_e)_*\bigl(a_{M,N,e}(k)\,dk\bigr).
\end{equation}
Its total variation is bounded by $\int|a_{M,N,e}|\,dk$.  Formula
\eqref{meas:cyclic-trace} shows that its $\chi_D$-moment is exactly the
orbit contribution with the Hecke insertion.

\begin{proposition}\label{meas:common-support}
Suppose $M,N$ belong to one retained component and have closed
$z$-support bounds $Z_M,Z_N\subset\mathfrak j$.  Their measure satisfies
\begin{equation}\label{meas:common-support-bound}
 \operatorname{supp}(\nu_{f_N,g_M})
 \subset
 \bigcup_{e\in Z_M\cap Z_N}B_{\Ghat.e}.
\end{equation}
\end{proposition}

\begin{proof}
Sum \eqref{meas:orbit-measure} over the finitely many stable nilpotent
orbits in the common support.  Lemma~\ref{meas:supported-convergence}
gives additivity, and Lemma~\ref{meas:hom-functions} identifies all its
moments.  Orbits outside the common support give zero differential Hom;
permuted orbits give zero trace.  The support bound follows from
\eqref{meas:parameter-map}.  Uniqueness follows from the density of
characters in $C(B)$ in Lemma~\ref{desc:compactsets}.
\end{proof}

By Theorem~\ref{tr:spanning}, every pair of functions is a finite linear
combination of the pairs just considered.  Each vector uses only finitely
many retained fixed components, because the categorical trace is the
direct sum on such components.  Cross-component Homs vanish.  Adding the
finite measures above therefore proves existence in
Theorem~\ref{meas:detector} for every $f,g$, and character density proves
uniqueness independent of the chosen expansion.  The construction is
bilinear.  Proposition~\ref{meas:common-support} already proves the forward
implication in \eqref{meas:criterion}.

\subsection{Nondegeneracy of the top-orbit contribution}

The decisive test concerns total mass. Canceling a common nonvanishing
density will isolate a simple coefficient, so injectivity of the
parameter map $\beta_e$ is unnecessary.

We must also show that a nonzero quotient class cannot be hidden by
cancellation of its measure.  Fix a stable orbit $O=J.e$.  Let $P,Q$ be
finite-dimensional equivariant $D_e$-modules, with compatible cyclic
structures, and take the compact orbit-closure lifts of
Proposition~\ref{orb:weil-lifts}.  We put the test coefficient $Q$ first
and the input coefficient $P$ second in a Hom.

On an open subset where $O$ is closed and the proper boundary of its
closure has been removed, both lifts are already supported on $O$.
Consequently local cohomology with that support is the actual Hom
complex.  There is no additional series of inverse normal powers to put
into its Euler character.  Proposition~\ref{orb:hom-calculation} gives
\begin{equation}\label{meas:top-density}
 a_{Q,P,e}(k)=p_e(k)\,
     \operatorname{tr}\bigl(k\tau_e;\Hom_{D_e}(Q,P)\bigr),
 \qquad k\in L_{e,c}.
\end{equation}
The action on the Hom uses the inverse structure on the test as in
\eqref{meas:hom-action}.  The factor $p_e$ is the graded Euler character of
\begin{equation}\label{meas:universal-density-factors}
 \C[\ker e]\otimes
 \Lambda^\bullet(N_e[-1])\otimes
 \Lambda^\bullet(\Lie(U_e)^*[-1]).
\end{equation}
The identification of constant Clifford super-intertwiners with ordinary
intertwiners, and its compatibility with lifted parity and tensor twists,
are part of Proposition~\ref{orb:hom-calculation}.

The invariant-base Koszul complexes used to obtain actual compact lifts
do not alter \eqref{meas:top-density} except by constant nonzero factors.
Indeed, expand their finite exterior terms by additivity in each input.
Their equation spaces are $J$-invariant, so these factors are independent
of $k$; their Frobenius eigenvalues and inverse eigenvalues are never one.
The exterior traces are therefore nonzero, as in
Lemma~\ref{tr:invariant-support}.  We rescale the trace classes by these
constants.  This permits computation first on the uncompleted orbit model
while retaining actual compact tests throughout.

\begin{lemma}[Cancellation of the universal density]
\label{meas:cancel-density}
The function $p_e$ is continuous and nowhere zero on $L_{e,c}$.  Its
reciprocal is a uniform limit of finite complex linear combinations of
twisted characters of finite-dimensional $L_e$-representations with
compatible transporter structures.  Such a twisted character can be
inserted in \eqref{meas:top-density} by tensoring the test coefficient with
the dual representation, in lifted degree zero.
\end{lemma}

\begin{proof}
Decompose each of the finite-dimensional spaces in
\eqref{meas:universal-density-factors} by lifted degree.  The symmetric
algebra contributes inverse determinants, and each exterior algebra
contributes determinants.  Thus $p_e$ is a finite product of terms
\begin{equation}\label{meas:determinant-factor}
 \det(1-\varepsilon kA;W)^{\eta},
 \qquad \varepsilon\in\{1,-1\},\quad \eta\in\{1,-1\},
\end{equation}
where $W$ is a finite-dimensional $L_e$-representation with transporter
$A$, and all eigenvalues of $kA$ have modulus strictly less than one,
uniformly in $k$.  The sign $\varepsilon$ incorporates the specified
lifted-degree shift.  The strict inequality follows from the weight
bounds: all factors outside the constant coefficient Hom have negative
absolute weight.  Each determinant in \eqref{meas:determinant-factor} is
therefore nonzero, proving the first assertion.

The identities
\[
 \det(1-A;W)=\sum_{j=0}^{\dim W}(-1)^j
                         \operatorname{tr}(A;\Lambda^jW),
 \qquad
 \det(1-A;W)^{-1}=\sum_{n\geq0}
                         \operatorname{tr}(A;\Sym^nW)
\]
apply also with the sign $\varepsilon$.  The second series converges
uniformly here: its $n$th term is bounded by
$\dim(\Sym^nW)\rho^n$ for a fixed $\rho<1$.  Finite products of these
identities express $1/p_e$ as the required uniform limit.  Every finite
term is the twisted character of a finite tensor product of symmetric and
exterior powers of the allowed representation data; scalar factors may
be put into its coefficient or cyclic structure.

Finally,
\[
 \Hom_{D_e}(Q\otimes R^*,P)
 \simeq R\otimes\Hom_{D_e}(Q,P)
\]
with the induced cyclic actions.  Under the parity identification already
used for \eqref{meas:top-density}, its character multiplies that density
by the twisted character of $R$.  The extension construction in
Proposition~\ref{orb:weil-lifts} applies to $Q\otimes R^*$ as well.  Each
finite approximant is consequently realized by a finite linear
combination of permitted compact test classes.
\end{proof}

\begin{lemma}[Top quotient detection]\label{meas:top-detection}
Every nonzero finite linear combination of fixed-simple classes in the
$O$-quotient is detected by a compact test supported on $\overline O$:
its top-orbit measure has nonzero total mass on $B_{\Ghat.e}$.
\end{lemma}

\begin{proof}
Write the quotient class as $u=\sum_P c_P[P]$, with finitely many nonzero
coefficients.  Suppose its top mass were zero against every permitted test
lift.  Fix a test coefficient $Q$.  By tensoring $Q$ with the duals of the
representations in Lemma~\ref{meas:cancel-density}, the assumption gives
\[
 \int_{L_{e,c}}p_e(k)h(k)
    \sum_P c_P\operatorname{tr}
          \bigl(k\tau_e;\Hom_{D_e}(Q,P)\bigr)\,dk=0
\]
for every finite character combination $h$ used there.  Letting these
$h$ converge uniformly to $1/p_e$ gives
\begin{equation}\label{meas:unweighted-coefficient-pairing}
 0=\sum_P c_P\operatorname{tr}
            \bigl(\tau_e;\Hom_{D_e,L_e}(Q,P)\bigr).
\end{equation}
This is legitimate integration against a fixed continuous function, not
a limit of automorphic functions.

The equivariant Clifford coefficient category is semisimple.  Choose
$Q$ to be any fixed simple appearing in $u$, with the matching structure.
For $P\not\simeq Q$ its Hom is zero; for $P\simeq Q$ it is the scalar
endomorphism line and the action is the identity.  Thus
\eqref{meas:unweighted-coefficient-pairing} says $c_Q=0$.  Repeating this
contradicts $u\ne0$.  Some actual finite test must therefore have nonzero
top mass.  The argument requires no injectivity of
$\beta_e:L_{e,c}\to B_{\Ghat.e}$: the total mass alone suffices.
\end{proof}

\subsection{Excluding boundary cancellation}

\begin{lemma}\label{meas:boundary-rank}
If $e'\in\overline{J.e}\setminus J.e$, then
$\dim(\Ghat.e')<\dim(\Ghat.e)$.  In particular the two ambient nilpotent
orbits, and their compact sets, are distinct.
\end{lemma}

\begin{proof}
Since $J$ is reductive in characteristic zero, choose a $J$-stable
complement $W$ in $\ghat=\mathfrak j\oplus W$.  The adjoint rank on
$\mathfrak j$ is the dimension of the $J$-orbit and strictly drops on a
proper boundary orbit.  The rank on $W$ cannot increase under
specialization, because the condition that a matrix have rank at most a
given integer is closed.  Hence
\[
 \rk(\operatorname{ad}(e')|\ghat)
 <\rk(\operatorname{ad}(e)|\ghat).
\]
These ranks are the ambient orbit dimensions.  The same argument works
for a disconnected $J$, whose orbits are finite unions of orbits of its
identity component.  Disjointness of the compact sets follows from
Lemma~\ref{desc:compactsets}.
\end{proof}

We finish the converse in Theorem~\ref{meas:detector}. Suppose
$f\notin\cF_Y\cA_{\C}$. By Theorem~\ref{tr:spanning}, choose a
retained fixed component on which its trace class does not lie in the
supported trace for $Y\cap\mathfrak j$. Starting with this closed subset,
filter the nilpotent cone in $\mathfrak j$ by closed invariant subsets,
adding one Frobenius orbit of $J$-orbits at a time, compatibly with
closure. Write $T_0\subset T_1\subset\cdots\subset T_m$ for the resulting
supported trace subspaces, with $T_0$ the trace for $Y\cap\mathfrak j$.
Let $i$ be the least index for which the
component class lies in $T_i$. Then $i>0$, and its image
$u\in T_i/T_{i-1}$ is nonzero. A nontrivially permuted stratum has zero
trace, so this step adds a single Frobenius-stable orbit $O=J.e$ outside
$Y\cap\mathfrak j$.

Express $u$ as a finite combination of the fixed-simple classes in the
$O$-quotient, and choose their compact orbit-closure lifts as in the proof
of Theorem~\ref{tr:spanning}. Subtract their sum from the component
class. The remainder belongs to $T_{i-1}$ and, by the same theorem,
is represented by a finite combination of compact Weil classes with
that earlier support. Apply
Lemma~\ref{meas:top-detection} to $u$ and these lifts to obtain a compact
test supported on $\overline O$ with nonzero top mass. The earlier
support bound defining $T_{i-1}$ meets $\overline O$ only in
$\overline O\setminus O$. Proposition~\ref{meas:common-support} and
Lemma~\ref{meas:boundary-rank} therefore show that the remainder,
and the proper-boundary contributions of the chosen lifts, have zero
mass on $B_{\Ghat.e}$. This also excludes contributions from incomparable
$J$-orbits in the same ambient nilpotent orbit. Other components give
zero Hom. Therefore
\[
 \nu_{f,g}(B_{\Ghat.e})\ne0
\]
for the resulting genuine finitely supported function test $g$.
Since $e\notin Y$, the invariant set $Y$ does not contain its ambient
orbit, and $B_{\Ghat.e}$ is disjoint from $B_Y$.  The asserted support
condition fails.  This proves the converse and completes the proof of
Theorem~\ref{meas:detector}.

No step used a rational point of $X$, a cuspidality condition, or
Hecke-finiteness.  Empty and full $Y$ are included, and the closed point
$v$ may have any positive degree.
